\chapter{Appendix: Algebra and Probability Reference}

This appendix is a compact, self-contained refresher for the two mathematical
languages the rest of the book speaks: the \emph{algebra} of groups, rings, and
fields (the stage on which RSA, Diffie--Hellman, and elliptic curves live), and
the \emph{probability} of uniform sampling, entropy, and collisions (the stage on
which keys, nonces, and hash security live). Nothing here is new if you read the
core chapters; it is gathered in one place so you can look it up without hunting.

The treatment is deliberately terse but worked. Every abstract definition is
followed by a tiny concrete object you can hold in your head, and every theorem
you actually use in cryptography is stated precisely with a one-line reason it is
true. Where a proof would take a page, I point to it and give the intuition
instead.

\begin{marginnotebox}
If you want the full development with all proofs, see the companion volumes:
\emph{Probability From the Ground Up} for distributions, entropy, and the
birthday/collision machinery, and \emph{Algorithms From the Ground Up} for the
cost of modular exponentiation, GCD, and polynomial multiplication. This appendix
gives you the statements; those books give you the derivations. Cross-references
appear in \texttt{Connection} boxes throughout.
\end{marginnotebox}

\chapterline

\section{The big picture: why cryptographers care about algebra}
\label{ch41:bigpicture}

Almost every public-key scheme has the same shape. You have a set with an
operation. Composing elements is cheap; \emph{undoing} the composition is (we
hope) hard. Encryption walks forward along the easy direction; security rests on
the adversary being stuck trying to walk backward.

\begin{itemize}
  \item In RSA the set is $\Zns{N}$ under multiplication; ``forward'' is raising
        to the public exponent $e$, ``backward'' is taking the $e$-th root, which
        is tied to factoring $N$.
  \item In Diffie--Hellman the set is a cyclic group $\G$; ``forward'' is
        exponentiation $g \mapsto g^x$, ``backward'' is the discrete logarithm.
  \item In elliptic-curve crypto the set is the points on a curve under a
        geometric addition law; ``forward'' is scalar multiplication $P \mapsto
        kP$, ``backward'' is again a discrete log.
\end{itemize}

To even \emph{state} these schemes you need the vocabulary of groups (for the
multiplicative and additive structures), rings (for arithmetic mod $N$, where you
add \emph{and} multiply but cannot always divide), and fields (for AES's byte
arithmetic and for the coordinate fields of elliptic curves). That is the tour
below.

\begin{intuitionbox}
A useful mental hierarchy: a \textbf{group} gives you one good operation (you can
combine and undo). A \textbf{ring} gives you two operations (add and multiply) but
division may fail. A \textbf{field} is a ring where division by anything nonzero
works --- it behaves like ordinary fractions. Each step up adds power and removes
freedom: there are many groups, fewer rings worth naming, and the finite fields
are pinned down completely (one for each prime power).
Figure~\ref{ch41:fig-hierarchy} shows the containment.
\end{intuitionbox}

\begin{figure}[H]
\centering
\begin{tikzpicture}[every node/.style={font=\small}]
  \draw[rounded corners=10pt, fill=sourceblue!10, draw=sourceblue!70]
        (-3.4,-2.2) rectangle (3.4,1.7);
  \node[anchor=north east, sourceblue!70!black] at (3.3,1.6)
        {\textbf{groups}\;(one operation)};
  \draw[rounded corners=10pt, fill=quotientorange!12, draw=quotientorange!70]
        (-2.7,-1.9) rectangle (2.7,0.9);
  \node[anchor=north east, quotientorange!60!black] at (2.6,0.8)
        {\textbf{rings}\;(two operations)};
  \draw[rounded corners=10pt, fill=targetgreen!14, draw=targetgreen!70]
        (-1.9,-1.6) rectangle (1.9,0.1);
  \node[anchor=north, targetgreen!50!black] at (0,0.0)
        {\textbf{fields}\;(division works)};
  \node at (0,-1.0) {$\F_p,\ \F_{2^8}$};
  \node[quotientorange!55!black] at (0,1.25) {$\Zn{n}$ ($n$ composite)};
  \node[sourceblue!60!black] at (0,-2.55) {$\Zns{n}$ under $\times$,\quad $\Zn{n}$ under $+$};
\end{tikzpicture}
\caption{The algebraic hierarchy. Each inner box has more structure: a field is a
special ring, and a ring's additive (or a field's multiplicative) part is a group.
Cryptography lives at every level.}
\label{ch41:fig-hierarchy}
\end{figure}

\section{Groups}
\label{ch41:groups}

\begin{definition}[Group]
\label{ch41:def-group}
A \emph{group} is a set $G$ with a binary operation $\cdot$ satisfying:
\begin{enumerate}
  \item \textbf{Closure:} $a\cdot b \in G$ for all $a,b\in G$.
  \item \textbf{Associativity:} $(a\cdot b)\cdot c = a\cdot(b\cdot c)$.
  \item \textbf{Identity:} there is $e\in G$ with $e\cdot a = a\cdot e = a$.
  \item \textbf{Inverses:} every $a$ has some $a^{-1}$ with $a\cdot a^{-1}=e$.
\end{enumerate}
If additionally $a\cdot b = b\cdot a$ for all $a,b$, the group is
\emph{abelian} (commutative). Nearly every group in cryptography is abelian.
\end{definition}

Two notations coexist. In \emph{multiplicative} notation the operation is
$\cdot$, the identity is $1$, and repeated operation is a power $a^k$. In
\emph{additive} notation the operation is $+$, the identity is $0$, and repeated
operation is a multiple $k\cdot a$. They describe the same idea; elliptic curves
use additive notation, $\Zns{N}$ uses multiplicative.

\begin{example}[The workhorse groups]
\label{ch41:ex-workhorse}
\leavevmode
\begin{itemize}
  \item $\Zn{n}$ under addition: $\{0,1,\dots,n-1\}$ with addition mod $n$.
        Identity $0$; inverse of $a$ is $n-a$. Always a group, for every $n$.
  \item $\Zns{n}$ under multiplication: the residues mod $n$ that are
        \emph{coprime} to $n$, with multiplication mod $n$. Identity $1$. The
        coprimality requirement is exactly what guarantees inverses exist (see
        \S\ref{ch41:rings}). Its size is $\varphi(n)$, Euler's totient.
\end{itemize}
\end{example}

Take $n=10$. Then $\Zns{10}=\{1,3,7,9\}$, since those are the residues sharing no
factor with $10$. Check that it is closed: $3\cdot7 = 21 \equiv 1$, $3\cdot3 = 9$,
$7\cdot9 = 63 \equiv 3$, all land back in the set. Here $\varphi(10)=4$.

\subsection{Order, generators, and cyclic groups}
\label{ch41:order}

\begin{definition}[Order of an element]
\label{ch41:def-order}
The \emph{order} of $a\in G$, written $\ord(a)$, is the smallest positive integer
$k$ with $a^k = e$ (or $k\cdot a = 0$ additively). The \emph{order of the group}
$\abs{G}$ is its number of elements.
\end{definition}

\begin{example}[Orders in $\Zns{10}$]
\label{ch41:ex-order-10}
Powers of $3$ mod $10$: $3^1=3,\ 3^2=9,\ 3^3=27\equiv7,\ 3^4=81\equiv1$. So
$\ord(3)=4$. Because $4=\abs{\Zns{10}}$, the powers of $3$ sweep out the whole
group: $3$ is a \emph{generator}.
\end{example}

\begin{definition}[Cyclic group, generator]
\label{ch41:def-cyclic}
$G$ is \emph{cyclic} if some $g\in G$ has $\ord(g)=\abs{G}$; then every element is
a power of $g$, written $G=\langle g\rangle$, and $g$ is a \emph{generator} (or
\emph{primitive root}).
\end{definition}

\begin{focusbox}
Cyclic groups are the beating heart of discrete-log cryptography. A generator $g$
gives a bijection between exponents $\{0,1,\dots,\abs{G}-1\}$ and group elements:
$x \mapsto g^x$. Diffie--Hellman is exactly the statement that this map is easy
forward and (in the right group) hard to invert.
\end{focusbox}

\begin{theorem}[Lagrange's theorem]
\label{ch41:thm-lagrange}
In a finite group $G$, the order of every element divides $\abs{G}$. Equivalently
$a^{\abs{G}} = e$ for all $a\in G$.
\end{theorem}

The reason is a counting trick: an element of order $k$ generates a subgroup of
exactly $k$ elements, and a finite group splits into equal-sized translated copies
(\emph{cosets}) of any subgroup, so $k$ must divide $\abs{G}$. Two cryptographic
payoffs fall straight out.

\begin{corollary}[Euler and Fermat]
\label{ch41:cor-euler}
For $\gcd(a,n)=1$, $a^{\varphi(n)} \equiv 1 \pmod n$ (Euler). When $n=p$ is prime
this is $a^{p-1}\equiv 1 \pmod p$ (Fermat's little theorem).
\end{corollary}

These two congruences are why RSA decryption works and why so many primality and
inverse computations terminate where they do.

\begin{connectionbox}
The cost of computing $g^x$ by repeated squaring --- $O(\log x)$ multiplications
--- is derived in \emph{Algorithms From the Ground Up} (fast exponentiation). The
number-theoretic facts that $\varphi$ is multiplicative and that $\Zns{p}$ is
cyclic are proved in the companion \emph{Appendix: Number Theory Reference}.
\end{connectionbox}

\subsection{Subgroups and the structure of $\Zns{p}$}

A \emph{subgroup} is a subset that is itself a group under the same operation. By
Lagrange its size divides $\abs{G}$. For a prime $p$, the group $\Zns{p}$ is
cyclic of order $p-1$, and its subgroups are exactly one for each divisor of
$p-1$. Cryptographers exploit this: to get a group of large \emph{prime} order $q$
(which kills small-subgroup attacks), pick a ``safe prime'' $p=2q+1$ and work in
the order-$q$ subgroup of $\Zns{p}$.

\begin{trapbox}
Working in $\Zns{p}$ \emph{as a whole} is dangerous when $p-1$ has small prime
factors: the Pohlig--Hellman attack solves the discrete-log piece by piece across
those small subgroups and recombines via the Chinese Remainder Theorem, turning a
``hard'' log into many easy ones. Restricting to a prime-order subgroup removes the
small factors. This is also why an attacker who can confine you to a tiny subgroup
(a \emph{small-subgroup confinement} attack) can leak bits of your secret ---
always validate that a received element has the expected large order.
\end{trapbox}

\section{Rings: where you can add and multiply}
\label{ch41:rings}

\begin{definition}[Commutative ring with unity]
\label{ch41:def-ring}
A \emph{ring} $R$ is a set with two operations, $+$ and $\cdot$, such that
$(R,+)$ is an abelian group with identity $0$, multiplication is associative and
has an identity $1$, and the \emph{distributive law} $a(b+c)=ab+ac$ holds. If
multiplication also commutes, the ring is \emph{commutative}. All rings in this
book are commutative with unity.
\end{definition}

The prototype is $\Zn{n}=\{0,1,\dots,n-1\}$ with addition and multiplication mod
$n$. You can always add, subtract, and multiply. What you cannot always do is
\emph{divide}.

\begin{definition}[Unit, zero divisor]
\label{ch41:def-unit}
An element $a$ is a \emph{unit} if it has a multiplicative inverse $a^{-1}$. It is
a \emph{zero divisor} if $a\neq 0$ yet $ab=0$ for some $b\neq 0$.
\end{definition}

In $\Zn{n}$ the element $a$ is a unit exactly when $\gcd(a,n)=1$, and the units
form the group $\Zns{n}$ of \S\ref{ch41:groups}. The non-units (other than $0$)
are zero divisors.

\begin{example}[Why $4$ has no inverse mod $10$]
\label{ch41:ex-zerodiv}
$\gcd(4,10)=2\neq1$, so $4$ is not a unit. Indeed $4\cdot 5 = 20 \equiv 0 \pmod
{10}$, so $4$ and $5$ are zero divisors: multiplying by $4$ collapses distinct
inputs to the same output, and a collapsing map cannot be undone. Contrast
$\gcd(3,10)=1$: the extended Euclidean algorithm gives $3\cdot 7 = 21 \equiv 1$,
so $3^{-1}\equiv 7$.
\end{example}

\begin{whybox}
\textbf{Why coprimality is exactly the condition for invertibility.} If
$\gcd(a,n)=1$, Bezout's identity gives integers $u,v$ with $au+nv=1$; reducing mod
$n$ leaves $au\equiv 1$, so $u=a^{-1}$. Conversely, if $\gcd(a,n)=d>1$, then $a$
times anything is a multiple of $d$, so it can never equal $1\bmod n$ unless $d\mid
1$. The extended Euclidean algorithm \emph{computes} the $u,v$, which is how
implementations find modular inverses (and RSA private exponents) in $O(\log n)$
steps.
\end{whybox}

\begin{connectionbox}
The extended Euclidean algorithm and its running time are in \emph{Algorithms From
the Ground Up}. The fact that $\Zn{n}$ splits as a product $\Zn{p_1^{e_1}} \times
\cdots$ via the Chinese Remainder Theorem --- the engine behind fast RSA
decryption --- is in the number-theory appendix.
\end{connectionbox}

\section{Fields}
\label{ch41:fields}

\begin{definition}[Field]
\label{ch41:def-field}
A \emph{field} is a commutative ring in which every nonzero element is a unit. So
you can add, subtract, multiply, and divide by anything nonzero --- arithmetic
behaves like the rationals or reals.
\end{definition}

The single most useful fact:

\begin{theorem}[$\Zn{p}$ is a field iff $p$ is prime]
\label{ch41:thm-fp}
$\Zn{n}$ is a field if and only if $n$ is prime. When $p$ is prime we write the
field as $\F_p$.
\end{theorem}

If $p$ is prime, every nonzero residue is coprime to $p$, hence a unit (by
\S\ref{ch41:rings}); so every nonzero element is invertible and $\Zn{p}$ is a
field. If $n$ is composite, $n=ab$ with $1<a,b<n$ gives zero divisors $a,b$, and a
zero divisor can never be a unit --- so $\Zn{n}$ is not a field.

\begin{example}[$\F_5$ by hand]
\label{ch41:ex-f5}
In $\F_5=\{0,1,2,3,4\}$ every nonzero element has an inverse: $1^{-1}=1$,
$2^{-1}=3$ (since $2\cdot3=6\equiv1$), $3^{-1}=2$, $4^{-1}=4$ (since $4\cdot4=16
\equiv1$). So $3/4 = 3\cdot 4^{-1} = 3\cdot 4 = 12 \equiv 2$. Division just works.
\end{example}

\subsection{Finite fields of prime-power order}
\label{ch41:gf}

Prime fields $\F_p$ are not the only finite fields. There is exactly one finite
field of size $q$ for every prime power $q=p^m$ (and none of any other size), all
called $\F_q$ or $\mathrm{GF}(q)$. The case cryptographers meet constantly is
$\F_{2^m}$ --- AES lives in $\F_{2^8}$ and many MACs and codes live in
$\F_{2^{128}}$.

You build $\F_{p^m}$ not from integers mod $p^m$ (that ring has zero divisors and
is \emph{not} a field) but from \emph{polynomials} over $\F_p$ reduced modulo an
\emph{irreducible} polynomial of degree $m$ --- the polynomial analogue of taking
integers mod a prime.

\begin{definition}[Polynomial ring and irreducibility]
\label{ch41:def-poly}
$\F_p[x]$ is the ring of polynomials with coefficients in $\F_p$, added and
multiplied as usual but with coefficient arithmetic done mod $p$. A nonconstant
$f(x)$ is \emph{irreducible} over $\F_p$ if it cannot be written as a product of
two lower-degree polynomials in $\F_p[x]$ --- the polynomial version of ``prime.''
\end{definition}

\begin{theorem}[Construction of $\F_{p^m}$]
\label{ch41:thm-gf}
If $f(x)\in\F_p[x]$ is irreducible of degree $m$, then $\F_p[x]/(f(x))$ --- the
polynomials of degree $<m$, with multiplication reduced mod $f$ --- is a field of
$p^m$ elements. Different choices of $f$ give isomorphic (``relabeled-identical'')
fields.
\end{theorem}

The reason it is a field exactly parallels the integer case: irreducibility of $f$
plays the role of primality of $p$, so the only polynomial sharing a factor with
$f$ is a multiple of $f$ (which is $0$ in the quotient), and the extended Euclidean
algorithm \emph{for polynomials} produces inverses.

\subsection{Worked example: arithmetic in $\F_{2^3}$}
\label{ch41:ex-gf8}

Take $p=2$, $m=3$, and the irreducible polynomial $f(x)=x^3+x+1$. Elements are the
$8$ polynomials of degree $<3$ over $\F_2$; written as bit-triples they are
$000,001,\dots,111$. Coefficient arithmetic is mod $2$, so \emph{addition is XOR}.

\textbf{Addition.} $(x^2+1)+(x+1) = x^2+x = $ in bits $101 \oplus 011 = 110$. No
carries, no reduction --- adding in $\F_{2^m}$ is just XOR of the bit-strings.

\textbf{Multiplication.} Compute $(x^2+1)\cdot(x+1)$:
\[
(x^2+1)(x+1) = x^3 + x^2 + x + 1.
\]
This has degree $3$, so reduce mod $f$. Since $x^3+x+1\equiv 0$, we have
$x^3 \equiv x+1$. Substitute:
\[
x^3 + x^2 + x + 1 \equiv (x+1) + x^2 + x + 1 = x^2 + 2x + 2 \equiv x^2 \pmod{2},
\]
because $2x$ and $2$ vanish mod $2$. So the product is $x^2$, i.e.\ bits $100$.

\begin{intuitionbox}
The reduction step ``$x^3\equiv x+1$'' is the polynomial twin of ``$10 \equiv 0$
mod $10$.'' Whenever a product overflows degree $m$, you fold the top term back in
using the relation the irreducible polynomial gives you. In $\F_{2^8}$ (AES) the
relation is $x^8 \equiv x^4+x^3+x+1$, which is the famous reduction by the byte
\texttt{0x11B}.
\end{intuitionbox}

\begin{implementbox}
Multiplication in $\F_{2^8}$ with AES's reduction polynomial \texttt{0x11B}, the
``Russian peasant'' way: shift-and-XOR, folding in the modulus whenever bit~$8$
appears.

\begin{lstlisting}[language=Python]
def gf_mul(a, b):
    # Multiply two bytes in GF(2^8) with AES modulus x^8+x^4+x^3+x+1 (0x11B).
    # LEARNING ONLY: not constant-time -- the data-dependent branch below
    # leaks the operands through timing. Never use for real crypto.
    p = 0
    for _ in range(8):
        if b & 1:            # add current shifted a (XOR = add in GF(2))
            p ^= a
        b >>= 1
        carry = a & 0x80     # will a overflow past bit 7?
        a <<= 1
        if carry:            # fold in the modulus: subtract x^8+x^4+x^3+x+1
            a ^= 0x11B
    return p & 0xFF

assert gf_mul(0x53, 0xCA) == 0x01   # 0x53 and 0xCA are inverses in AES's field
\end{lstlisting}
\textbf{Warning:} this is toy code to show the algebra. A production
implementation uses constant-time table lookups or carry-less multiplication
(\texttt{PCLMULQDQ}) to avoid the timing side channel flagged above.
\end{implementbox}

\begin{applicationbox}
\textbf{Where finite fields show up.} AES does every \texttt{MixColumns} step as
matrix multiplication over $\F_{2^8}$. AES-GCM's authentication tag is a polynomial
evaluated over $\F_{2^{128}}$ (a universal hash). Reed--Solomon codes (CDs, QR
codes, deep-space links) and Shamir secret sharing both interpolate polynomials
over a finite field. Elliptic curves are defined by an equation whose coordinates
live in $\F_p$ or $\F_{2^m}$. If you understand $\F_{2^3}$ above, you understand
the algebra under all of these --- only the sizes change.
\end{applicationbox}

\section{A pocket card of algebraic facts}
\label{ch41:card}

\begin{center}
\renewcommand{\arraystretch}{1.25}
\begin{tabular}{@{}p{0.40\textwidth}p{0.54\textwidth}@{}}
\toprule
\textbf{Fact} & \textbf{Where it is used} \\
\midrule
$a^{-1} \bmod n$ exists iff $\gcd(a,n)=1$ & RSA key gen, modular division \\
$a^{\varphi(n)}\equiv 1 \pmod n$ (Euler) & RSA correctness \\
$a^{p-1}\equiv 1 \pmod p$ (Fermat) & primality tests, $a^{-1}=a^{p-2}$ \\
$\Zns{p}$ is cyclic of order $p-1$ & Diffie--Hellman, ElGamal \\
$\ord(g) \mid \abs{G}$ (Lagrange) & subgroup choice, Pohlig--Hellman \\
$\Zn{n}$ is a field iff $n$ prime & choosing curve/field parameters \\
$\F_{p^m}=\F_p[x]/(f)$, $f$ irreducible & AES, GCM, Reed--Solomon, sharing \\
Add in $\F_{2^m}$ is bitwise XOR & every byte-level cipher operation \\
\bottomrule
\end{tabular}
\end{center}

\section{Probability: the cryptographer's slice}
\label{ch41:prob}

Cryptography uses a narrow but crucial corner of probability: \emph{uniform}
sampling (good keys, good nonces), \emph{entropy} (how much real
unpredictability you have), and \emph{collision} bounds (when do random draws
repeat). We collect those here.

\subsection{Probability spaces and uniform distributions}
\label{ch41:uniform}

A finite \emph{probability space} is a finite set $\Omega$ of outcomes with weights
$\Pr[\omega]\ge 0$ summing to $1$. An \emph{event} is a subset $A\subseteq\Omega$
with $\Pr[A]=\sum_{\omega\in A}\Pr[\omega]$. A \emph{random variable} is a function
$X:\Omega\to\R$ (or into any set).

The distribution cryptography cares about most is the \emph{uniform} one:
$\Pr[\omega]=1/\abs{\Omega}$ for every outcome. Writing $X\getsr S$ means $X$ is
drawn uniformly from the set $S$. A $128$-bit key chosen uniformly from
$\{0,1\}^{128}$ is the canonical example: every one of the $2^{128}$ strings is
equally likely, and \emph{that} is what makes guessing it hopeless.

\begin{definition}[Statistical (total variation) distance]
\label{ch41:def-sd}
For distributions $P,Q$ over the same finite set $S$,
\[
\Delta(P,Q) \;=\; \tfrac12 \sum_{s\in S}\,\abs{\,P(s)-Q(s)\,}
            \;=\; \max_{A\subseteq S}\,\abs{\,P(A)-Q(A)\,}.
\]
It is $0$ when the distributions are identical and $1$ when their supports are
disjoint.
\end{definition}

Statistical distance is the yardstick for ``looks uniform.'' A PRG output is good
if its distribution is statistically (or, weaker but enough, computationally) close
to uniform; an attacker who cannot tell the two apart cannot exploit the
difference. The triangle inequality $\Delta(P,R)\le \Delta(P,Q)+\Delta(Q,R)$ lets
you chain these closeness arguments through a proof --- the ``hybrid argument'' you
meet in security reductions.

\begin{example}[A biased coin is detectable]
\label{ch41:ex-bias}
Let $P$ be a fair bit ($\Pr[0]=\Pr[1]=\tfrac12$) and $Q$ a coin with
$\Pr[1]=0.51$. Then $\Delta(P,Q)=\tfrac12(\abs{0.5-0.49}+\abs{0.5-0.51})=0.01$.
Small --- but over many samples a distinguisher that always guesses ``$Q$ when it
sees more $1$s'' gains an edge that grows, which is exactly how statistical tests
catch a broken RNG.
\end{example}

\subsection{Expectation, variance, and the union bound}
\label{ch41:expectation}

For a real random variable $X$, the \emph{expectation} is $\mathbb{E}[X] =
\sum_x x\cdot\Pr[X=x]$, and it is \emph{linear}: $\mathbb{E}[X+Y]=\mathbb{E}[X]
+\mathbb{E}[Y]$ \emph{even when $X,Y$ are dependent}. Linearity is the most
overpowered tool in the elementary toolkit --- the birthday count below uses it.

Two inequalities do most of the heavy lifting in security proofs.

\begin{theorem}[Union bound]
\label{ch41:thm-union}
For any events $A_1,\dots,A_k$, $\;\Pr[A_1\cup\cdots\cup A_k]\le\sum_{i}\Pr[A_i]$.
\end{theorem}

\begin{theorem}[Markov's inequality]
\label{ch41:thm-markov}
For a nonnegative $X$ and $t>0$, $\;\Pr[X\ge t]\le \mathbb{E}[X]/t$.
\end{theorem}

The union bound is how we say ``if no single bad event is likely, then \emph{some}
bad event happening is still unlikely'' --- e.g.\ bounding the chance that any of
$q$ encryption queries reuses a nonce. Markov turns an expectation into a tail
guarantee when that is all you have.

\subsection{Independence and conditional probability}
\label{ch41:indep}

Events $A,B$ are \emph{independent} if $\Pr[A\cap B]=\Pr[A]\Pr[B]$. The
\emph{conditional} probability is $\Pr[A\mid B]=\Pr[A\cap B]/\Pr[B]$, and Bayes's
rule rearranges it:
\[
\Pr[A\mid B] = \frac{\Pr[B\mid A]\,\Pr[A]}{\Pr[B]}.
\]
Independence is what lets a one-time-pad ciphertext reveal nothing: if the key bit
is uniform and independent of the message bit, the XOR is uniform regardless of the
message, so $\Pr[\text{message}\mid\text{ciphertext}]=\Pr[\text{message}]$ ---
perfect secrecy is exactly a statement about conditional probability.

\begin{connectionbox}
Full constructions of probability spaces, the proofs that expectation is linear and
that Markov/Chebyshev/Chernoff hold, and the law of large numbers (why statistical
tests on RNGs converge) are developed in \emph{Probability From the Ground Up}.
Here we only need the statements.
\end{connectionbox}

\subsection{Entropy: measuring unpredictability}
\label{ch41:entropy}

A key is only as strong as it is unpredictable, and ``unpredictable'' has a precise
measure. Let $X$ be a random variable taking value $x$ with probability $p_x$.

\begin{definition}[Shannon and min-entropy]
\label{ch41:def-entropy}
The \emph{Shannon entropy} (in bits) is
\[
H(X) = -\sum_x p_x \log_2 p_x,
\]
the average surprise. The \emph{min-entropy} is
\[
H_\infty(X) = -\log_2\Bigl(\max_x p_x\Bigr),
\]
governed by the single most likely value.
\end{definition}

\begin{trapbox}
For \emph{key material the quantity that matters is min-entropy, not Shannon
entropy.} Shannon entropy is an average; an attacker does not attack the average,
they attack the easiest case. If one password is chosen $50\%$ of the time and a
million others split the rest, the Shannon entropy looks healthy but
$H_\infty = -\log_2(0.5)=1$ bit: a single guess succeeds half the time. Always
size key-derivation and password requirements against $H_\infty$.
\end{trapbox}

\begin{example}[Entropy of a uniform key]
\label{ch41:ex-keyentropy}
For $X$ uniform on $\{0,1\}^{n}$, every value has $p_x = 2^{-n}$, so
$H(X)=H_\infty(X)=n$ bits --- the two notions coincide exactly when the
distribution is uniform. This is the formal sense in which ``a $128$-bit key has
$128$ bits of security against guessing.'' Any nonuniformity strictly lowers
$H_\infty$, which is why CSPRNG output must be statistically close to uniform
(\S\ref{ch41:uniform}).
\end{example}

\begin{applicationbox}
When you run \texttt{ssh-keygen} or a TLS library samples a session key, it pulls
bytes from the OS CSPRNG (\texttt{getrandom}, \texttt{/dev/urandom}), which is
seeded to have full min-entropy. Password managers target $\sim 128$ bits of
min-entropy by drawing characters uniformly from a known alphabet --- e.g.\
$20$ characters from a $90$-symbol set gives $20\log_2 90 \approx 130$ bits. The
moment a generator's output is biased (Debian's 2008 OpenSSL bug, which shrank the
key space to a few thousand values) the effective entropy collapses and keys become
guessable.
\end{applicationbox}

\subsection{The birthday bound: when do random draws collide?}
\label{ch41:birthday}

The single most-cited probability fact in cryptography. Draw $q$ values uniformly
and independently from a set of size $N$. How large must $q$ be before two draws
\emph{collide} (coincide)?

\begin{theorem}[Birthday bound]
\label{ch41:thm-birthday}
With $q$ uniform independent draws from $N$ values, the probability of at least one
collision satisfies
\[
\Pr[\text{collision}] \;\le\; \binom{q}{2}\frac{1}{N} \;=\; \frac{q(q-1)}{2N},
\]
and is well-approximated by $1-e^{-q^2/(2N)}$. The probability crosses $1/2$ around
$q \approx 1.177\sqrt{N}$, i.e.\ at roughly $\sqrt{N}$ draws
(Figure~\ref{ch41:fig-birthday}).
\end{theorem}

\begin{figure}[H]
\centering
\begin{tikzpicture}
\begin{axis}[
  width=0.82\textwidth, height=5.2cm,
  xlabel={draws $q$ (as a multiple of $\sqrt{N}$)},
  ylabel={$\Pr[\text{collision}]$},
  domain=0:3, samples=80,
  xmin=0, xmax=3, ymin=0, ymax=1.02,
  ytick={0,0.5,1}, xtick={0,1,1.177,2,3},
  xticklabels={$0$,$\sqrt N$,$1.18\sqrt N$,$2\sqrt N$,$3\sqrt N$},
  axis lines=left, every axis plot/.append style={thick},
  legend pos=south east, legend cell align=left,
]
\addplot[sourceblue] {1 - exp(-x^2/2)};
\addlegendentry{$1-e^{-q^2/(2N)}$}
\draw[densely dashed, gray] (axis cs:1.177,0) -- (axis cs:1.177,0.5)
      -- (axis cs:0,0.5);
\end{axis}
\end{tikzpicture}
\caption{Collision probability against the number of draws, with $q$ measured in
units of $\sqrt{N}$. The curve crosses one half at about $1.18\sqrt{N}$ draws,
regardless of how large $N$ is --- the square-root law in one picture.}
\label{ch41:fig-birthday}
\end{figure}

\begin{whybox}
\textbf{Why $\sqrt{N}$ and not $N$.} A collision needs a \emph{pair}, and there are
$\binom{q}{2}\approx q^2/2$ pairs among $q$ draws. Each specific pair coincides with
probability $1/N$. By linearity of expectation, the expected number of colliding
pairs is $\binom{q}{2}/N$; the union bound turns that into the upper bound above.
When the expected number of collisions reaches about $1$ --- which happens at
$q^2/2 \approx N$, i.e.\ $q\approx\sqrt{2N}$ --- collisions become likely. The
counterintuitive part is that you only need on the order of $\sqrt N$ draws, not
$N$, because the number of \emph{pairs} grows quadratically.
\end{whybox}

\begin{example}[Classic birthdays and a 64-bit tag]
\label{ch41:ex-birthday}
With $N=365$ days, $1.177\sqrt{365}\approx 22.5$ --- just $23$ people give a better
than even chance two share a birthday. Now scale up: a $64$-bit value has
$N=2^{64}$, so collisions become likely around $q\approx 2^{32}$ --- about four
billion draws. That is why a $64$-bit authentication tag or a $64$-bit block cipher
(such as 3DES) is risky at high volume, and why modern designs use $128$-bit tags
($\sqrt{N}=2^{64}$ draws, comfortably out of reach).
\end{example}

\begin{attackbox}
\textbf{Birthday attack on hash collisions.} To forge a document collision against
an $n$-bit hash, an attacker does \emph{not} search $2^n$ inputs. They generate
about $2^{n/2}$ variants of an innocuous message and $2^{n/2}$ variants of a
malicious one (whitespace, synonyms), hash them all, and look for any matching
pair. By the birthday bound a collision is likely after $\sim 2^{n/2}$ work.
This halves the effective security: MD5's $128$-bit digest gives only $\sim 2^{64}$
collision resistance (and is broken outright in practice), which is precisely why
collision-resistant hashes need an output \emph{twice} the target security level
--- SHA-256 for $128$-bit collision resistance.
\end{attackbox}

\begin{implementbox}
Watch the birthday bound appear empirically: draw random values and count how many
draws until the first collision, averaged over trials.

\begin{lstlisting}[language=Python]
import random, math

def first_collision(N, trials=2000):
    # Average number of uniform draws from [0, N) until the first repeat.
    # LEARNING ONLY: a demonstration, not a cryptographic RNG test.
    total = 0
    for _ in range(trials):
        seen = set()
        while True:
            x = random.randrange(N)
            if x in seen:
                break
            seen.add(x)
            total += 1
    return total / trials

for N in (365, 4096, 65536):
    print(N, round(first_collision(N)), "predicted", round(math.sqrt(math.pi*N/2)))
# Empirical count tracks sqrt(pi*N/2), the birthday-bound mean.
\end{lstlisting}
\textbf{Warning:} Python's \texttt{random} is a Mersenne Twister --- fine for this
illustration, never for keys or nonces. Use \texttt{secrets} or
\texttt{os.urandom} for anything security-relevant.
\end{implementbox}

\subsection{The negligible-function discipline}
\label{ch41:negl}

Security definitions need a notion of ``so small we do not care.'' A function
$\negl(\lambda)$ is \emph{negligible} if it shrinks faster than any inverse
polynomial: for every $c$, eventually $\negl(\lambda) < \lambda^{-c}$. Think
$2^{-\lambda}$ or $\lambda^{-\log\lambda}$. The whole apparatus of provable
security reads ``every $\PPT$ adversary wins with probability at most
$\negl(\lambda)$.'' Negligible functions are closed under addition and under
multiplication by polynomials, which is exactly what makes union-bound and hybrid
arguments compose without the error blowing up.

\begin{connectionbox}
The asymptotics behind ``negligible,'' and the formal definition of a $\PPT$
(probabilistic polynomial-time) machine, sit at the boundary of
\emph{Probability From the Ground Up} and \emph{Algorithms From the Ground Up}; the
computational-hardness chapter of this book builds the security games on top of
them.
\end{connectionbox}

\section{A pocket card of probability facts}
\label{ch41:probcard}

\begin{center}
\renewcommand{\arraystretch}{1.25}
\begin{tabular}{@{}p{0.42\textwidth}p{0.52\textwidth}@{}}
\toprule
\textbf{Fact} & \textbf{Where it is used} \\
\midrule
$X\getsr\{0,1\}^n$ has $H=H_\infty=n$ & key/nonce strength \\
Min-entropy $H_\infty=-\log_2\max_x p_x$ & passwords, KDF sizing \\
$\Delta(P,Q)=\tfrac12\sum\abs{P-Q}$ & PRG/indistinguishability proofs \\
Triangle ineq.\ for $\Delta$ & hybrid arguments \\
Union bound $\Pr[\cup A_i]\le\sum\Pr[A_i]$ & bounding many bad events \\
Linearity of $\mathbb{E}$ (always) & counting collisions \\
Collision likely at $q\approx\sqrt N$ & tag/block/hash sizing \\
$n$-bit hash $\Rightarrow 2^{n/2}$ collision work & choosing digest length \\
$\negl$: smaller than any $1/\mathrm{poly}$ & security definitions \\
\bottomrule
\end{tabular}
\end{center}

\section{How the two halves meet}
\label{ch41:synthesis}

The algebra gives you objects whose forward operation is easy and backward
operation is hard; the probability gives you the language to say \emph{how} hard,
and to size your parameters so that ``hard'' stays hard. A discrete-log group of
order $q$ resists generic attacks up to about $\sqrt{q}$ work --- a \emph{birthday
bound} again (Pollard's rho), which is why a $256$-bit elliptic curve targets
$128$-bit security. RSA needs a $\sim 3072$-bit modulus for the same $128$ bits
because factoring is sub-exponential, not square-root. In both cases you are
balancing an algebraic hardness assumption against a probabilistic attack cost.
Hold the two pocket cards side by side and most parameter choices in the book stop
being magic numbers and start being arithmetic.

\section{Exercises}
\label{ch41:exercises}

\begin{exercise}
List the elements of $\Zns{18}$ and verify $\abs{\Zns{18}}=\varphi(18)$. Find the
order of $5$ in this group, and decide whether $5$ is a generator.
\end{exercise}

\begin{exercise}
Using the extended Euclidean algorithm, compute $17^{-1} \bmod 43$. Check your
answer by multiplying. Then explain in one sentence why $\gcd$ being $1$ was
necessary.
\end{exercise}

\begin{exercise}
Show directly that $\Zn{12}$ is not a field by exhibiting two nonzero elements
whose product is $0$. Which elements of $\Zn{12}$ \emph{are} units?
\end{exercise}

\begin{exercise}
In $\F_{2^3}$ with modulus $f(x)=x^3+x+1$, compute $(x^2+x)\cdot(x+1)$ and reduce.
Then find the multiplicative inverse of $x$ (hint: what does $x\cdot x^2$ reduce
to?).
\end{exercise}

\begin{exercise}
The AES field $\F_{2^8}$ uses the reduction polynomial $x^8+x^4+x^3+x+1$. Verify
that multiplying the byte \texttt{0x02} by an element is ``shift left by one, and
XOR \texttt{0x1B} if the high bit was set.'' Compute $\texttt{0x02} \cdot
\texttt{0x80}$ and $\texttt{0x02}\cdot\texttt{0x57}$.
\end{exercise}

\begin{exercise}
A coin lands heads with probability $0.6$. Compute its Shannon entropy and its
min-entropy in bits, and explain which one you would use to budget guessing
resistance.
\end{exercise}

\begin{exercise}
You generate random $40$-bit nonces. Using the birthday approximation, roughly how
many nonces can you issue before the probability of a repeat reaches $1/2$? Restate
your answer as ``about $2^{?}$ nonces.''
\end{exercise}

\begin{exercise}
Prove that if $P,Q,R$ are distributions on a finite set, then $\Delta(P,R)\le
\Delta(P,Q)+\Delta(Q,R)$. (Use the definition $\Delta=\tfrac12\sum\abs{\cdot}$ and
the ordinary triangle inequality on each term.)
\end{exercise}

\begin{exercise}
An attacker wants a collision in a $96$-bit hash. About how much work (as a power
of two) does the birthday attack of \S\ref{ch41:birthday} require, and how many
bits of collision security does a $96$-bit digest therefore provide? Is that enough
for $128$-bit-security applications?
\end{exercise}

\begin{exercise}
Using the union bound, show that if each of $q$ independent encryptions draws a
nonce uniformly from $N$ values, the probability that \emph{any} two share a nonce
is at most $q^2/(2N)$. For $N=2^{96}$, how large can $q$ be while keeping this below
$2^{-32}$?
\end{exercise}

\begin{exercise}
Explain, using min-entropy, why a key derived from a password with $20$ bits of
min-entropy through any deterministic KDF still has at most $20$ bits of security,
no matter how strong the KDF is. What does this tell you about the role of a slow
KDF (such as Argon2)?
\end{exercise}
